\documentclass[11pt]{article}
\usepackage[margin=1in]{geometry}
\usepackage{graphicx}
\usepackage{url}
\usepackage{booktabs}
\usepackage{array}

\title{Master Project: Bridging Vision and Language---A Multi-Image AI Pipeline for Context-Aware NPC Storytelling}
\author{Hassan Coulibaly}

\begin{document}

\maketitle

\begin{abstract}
Modern interactive media frequently struggles with ludonarrative dissonance, where dynamically driven player actions conflict with static, pre-scripted game narratives. While recent advancements in Vision-Language Models (VLMs) have enabled narrative extraction from single, static images, these approaches lack the temporal and sequential context required to interpret complex gameplay. This project introduces a novel, multi-modal machine learning pipeline designed to generate context-aware, dynamic dialogue for Non-Playable Characters (NPCs) based on sequential gameplay screenshots. The proposed architecture separates visual perception from narrative generation to strictly ground the AI in gameplay reality. First, a VLM sequentially processes in-game images to extract key entities, actions, and environmental states, structuring this data into a chronological scene graph. Second, a Large Language Model (LLM) utilizes this structured data---constrained by established interactive storytelling principles derived from Adams~\cite{Adams2014}---to generate dramatically meaningful and credible NPC dialogue. By filtering raw visual data into a strict JSON framework before text generation, the pipeline reduces LLM hallucination and arbitrary content creation. The pipeline was evaluated on six \emph{Minecraft} runs stratified by context length (5, 10, and 15 minutes of gameplay, two runs each). Automated LLM-as-judge scoring across seven Adams-derived dimensions yielded an average Grounding score of 4.50/5, with all other dimensions (Coherence, Credibility, Repetition, Arbitrary Content, Emotional Richness, Engagement) averaging 5.00/5. Human raters largely agreed on Grounding (4.67), Repetition (5.00), and Arbitrary Content (5.00), but rated Emotional Richness (3.33) and Engagement (4.00) notably lower than the automated judge, suggesting that LLM-based judges may overestimate emotional quality relative to human perception. Results also indicate that longer context windows produce higher-quality narratives, with 15-minute runs receiving consistently higher scores than 5-minute runs from both evaluators.
\end{abstract}

\section{Introduction}

The landscape of interactive media and video games has evolved from simple, linear progression into expansive, open-world environments where player agency dictates the flow of events. As rendering technologies and physics engines have advanced to create highly realistic simulations, the systems governing in-game storytelling have often lagged behind. In most contemporary open-world games, Non-Playable Character (NPC) dialogue remains heavily pre-scripted. While branching dialogue trees offer the illusion of choice, they are ultimately finite and struggle to account for the emergent, unpredictable actions players take in a physics-driven world.

This master's project implements a prototype \emph{talking agent} pipeline: a batch of time-ordered gameplay screenshots is summarized per frame by a VLM, merged into a single structured \emph{scene graph}, and passed to an LLM that produces short NPC dialogue under explicit grounding constraints. The implementation uses a configurable Google Gemini model (default \texttt{gemini-2.5-flash}; overridable via \texttt{GEMINI\_MODEL} in \texttt{.env} or the \texttt{--model} flag) for both the vision and language calls, with extraction forced to \texttt{application/json} for machine-parseable state.

A key design goal is evaluation completeness. Prior work on VLM pipelines often evaluates only the perception stage. This project contributes a two-track evaluation: automated F1 for Stage~1, and an LLM-as-judge scoring method for Stage~3 whose dimensions are grounded in Adams's theory of game storytelling~\cite{Adams2014}. Adams defines a good game story as an account of dramatically meaningful events that is credible, coherent, without undue repetition, and without arbitrary content, and identifies length, characters, degree of realism, and emotional richness as key quality factors. These criteria directly motivate the seven scoring dimensions used in this project's dialogue evaluation. To study the effect of context length on narrative quality, the evaluation is stratified by gameplay duration: 5, 10, and 15 minutes of \emph{Minecraft} footage at a 5-second frame extraction interval, with two runs at each duration level.

\section{Background}

A prototype system for AI \& Interactive Storytelling in Cultural Heritage was developed for the Dunhuang Murals (specifically the Jataka stories) to move beyond passive viewing~\cite{11227644}. The system includes mural segmentation to explain individual elements, interactive knowledge graphs that trigger cinematic visual effects when clicked, and a creative reconstruction tool where users build their own comic-style narratives. An AI agent (built on the COZE platform and GPT models) acts as an approachable guide, providing real-time explanations and creative suggestions while drawing from a curated knowledge base to prevent ``hallucinations''. User testing showed high scores in attractiveness, novelty, and stimulation, suggesting the system helps users form deeper emotional and intellectual connections with heritage.

\section{System Architecture and Implementation}

\subsection{Design goals}

The design follows three principles aligned with the research abstract: (1)~\textbf{temporal grounding}---images are processed in a user-defined order (e.g., lexicographic filename order as a proxy for capture order); (2)~\textbf{structured mediation}---raw pixels are never passed directly to the dialogue model; only JSON derived from vision passes through; (3)~\textbf{constrained narration}---the final prompt instructs the model to avoid inventing entities, locations, or plot beats absent from the scene graph.

\subsection{Stage 1: Per-frame visual extraction}

For each frame index $i = 0,1,\ldots$, the pipeline sends the screenshot and a fixed instruction prompt to the VLM. The model must return a JSON object with fields including \texttt{location}, \texttt{time\_of\_day}, \texttt{player\_status}, \texttt{recent\_action}, \texttt{notable\_enemies\_or\_objects}, and \texttt{visible\_entities} (a list of short strings). Extraction uses a low temperature (0.2) to favor factual stability. Each result is wrapped with metadata: \texttt{frame\_index}, \texttt{source\_image}, and \texttt{observation}.

\subsection{Stage 2: Consolidation into a chronological scene graph}

Per-frame observations are merged into one JSON \emph{scene graph} suitable for dialogue generation. Two modes are implemented:

\begin{itemize}
  \item \textbf{LLM consolidation (default).} A second call to the same model merges the ordered frame list into a graph with \texttt{timeline\_summary}, \texttt{environment} (\texttt{location}, \texttt{time\_of\_day}, \texttt{conditions}), \texttt{entities} (with \texttt{id}, \texttt{name\_or\_description}, \texttt{type}, \texttt{first\_frame}, \texttt{last\_frame}), \texttt{events} (\texttt{frame\_index}, \texttt{description}, \texttt{involved\_entity\_ids}), and \texttt{player\_arc} (\texttt{status\_notes}, \texttt{notable\_actions}). The consolidator is instructed not to invent content absent from the observations.
  \item \textbf{Simple stitch (\texttt{--simple-consolidate}).} A deterministic merge (no extra API call) concatenates actions, collects entity strings, and builds a minimal graph; useful for cost-sensitive runs or ablations.
\end{itemize}

\subsection{Stage 3: NPC dialogue generation}

The LLM receives only the consolidated scene graph and a block of \emph{storytelling constraints} aligned with the Adams criteria: ground all claims in the JSON; target 80--120 words; avoid game-UI elements (health bars, menus, HUD text); express genuine emotion; and avoid repeating the same idea. Dialogue generation uses a higher temperature (0.7) for phrasing variety while the input remains fixed structured state.

\section{Experimental Setup}

\subsection{Dataset}

A 24-minute 21-second \emph{Minecraft} gameplay session was recorded and frames were extracted at a 5-second interval, yielding 292 total frames. Six structured runs were produced by varying context length:

\begin{itemize}
  \item \textbf{5-minute runs} (60 frames, 2 runs)
  \item \textbf{10-minute runs} (120 frames, 2 runs)
  \item \textbf{15-minute runs} (180 frames, 2 runs)
\end{itemize}

For each duration group, the first run executed the full three-stage pipeline (VLM extraction + LLM consolidation + dialogue generation). The second run reused the VLM-extracted frames from the first and ran only the LLM stages, testing output consistency given identical visual input.

\subsection{Codebase and Data Layout}

The repository is organized as follows (paths relative to project root):

\begin{verbatim}
talking_agent/
  src/
    main_pipeline.py          <- full 3-stage pipeline
    evaluator.py              <- Stage 1 F1 scoring
    dialogue_evaluator.py     <- Stage 3 LLM-as-judge scoring
    eda_visualize.py          <- EDA plots and figures
  data/
    test_screenshots/         <- input frames
    extracted_state/          <- *_frames.json, *_scene_graph.json
    generated_stories/        <- *_npc_dialogue.txt
    analysis/figures/         <- PNG exports from eda_visualize.py
  evaluation/
    ground_truth_labels.csv   <- binary frame labels for Stage 1
    dialogue_scores.csv       <- LLM-as-judge results
    dialogue_scores.json      <- same data, structured JSON
    human_scoring_template.csv <- human ratings
  .env.example
  requirements.txt
\end{verbatim}

\section{Evaluation}

Evaluation spans all three pipeline stages. Stage~1 is evaluated with automated binary classification metrics. Stage~3 is evaluated with an LLM-as-judge method and a human rating study. Stage~2 (scene graph consolidation) is assessed indirectly through the Grounding dimension of the Stage~3 judge: if consolidation omits or distorts key entities, the dialogue scores will reflect it.

\subsection{Evaluation framework: Adams's storytelling criteria}

The dialogue evaluation framework is derived from Adams~\cite{Adams2014}, who defines a good game story as an account of dramatically meaningful events satisfying five qualities: \textbf{credibility}, \textbf{coherence}, absence of \textbf{undue repetition}, absence of \textbf{arbitrary content}, and \textbf{dramatic meaning} (engagement). Adams also identifies \textbf{length}, character presence, \textbf{degree of realism}, and \textbf{emotional richness} as key story factors. These criteria map directly onto the seven scoring dimensions shown in Table~\ref{tab:dimensions}.

\begin{table}[h]
\centering
\caption{Dialogue scoring dimensions and their source in Adams~\cite{Adams2014}.}
\label{tab:dimensions}
\begin{tabular}{>{\bfseries}lp{7.5cm}}
\toprule
Dimension & Definition (1--5 scale) \\
\midrule
Grounding         & All concrete claims traceable to the scene graph; no invented entities or events. \\
Coherence         & Internally consistent, logically ordered, single narrative voice. \\
Credibility       & Believable within the game universe; tone and vocabulary fit an NPC companion. \\
Repetition        & Free from undue repetition of ideas or phrases (5 = no repetition). \\
Arbitrary Content & No irrelevant game-UI elements such as health bars, menus, or HUD text (5 = none). \\
Emotional Richness & Genuine emotion and human warmth, not a dry event log. \\
Engagement        & Dramatically meaningful and immersive as a campfire story. \\
\bottomrule
\end{tabular}
\end{table}

\subsection{Stage 1: VLM extraction accuracy (F1)}

The script \texttt{src/evaluator.py} reads \texttt{evaluation/ground\_truth\_labels.csv} with columns \texttt{Image\_ID}, \texttt{Condition\_Tested}, \texttt{Ground\_Truth}, and \texttt{VLM\_Prediction} (binary 0/1). It reports precision, recall, and F1 per condition and a macro-averaged F1. Conditions include binary checks such as \texttt{location\_detectable} and \texttt{action\_described}. Rows are populated with labels from the experimenter's rubric and pipeline outputs.

\subsection{Stage 3: Dialogue quality via LLM-as-judge}

The script \texttt{src/dialogue\_evaluator.py} implements an automated judge using the Adams dimensions above. For each run, it loads \texttt{*\_scene\_graph.json} and \texttt{*\_npc\_dialogue.txt}, builds a structured prompt that cites the Adams framework, and calls \texttt{gemini-2.5-flash} as the judge model, scoring all seven dimensions on a 1--5 scale. The judge also returns a list of \textbf{hallucinated entities}---proper nouns present in the dialogue but absent from the scene graph.

\subsubsection*{Automated results}

Table~\ref{tab:auto_scores} reports scores for all six Minecraft runs.

\begin{table}[h]
\centering
\caption{LLM-as-judge scores for all six Minecraft runs (scale: 1--5). G=Grounding, C=Coherence, Cr=Credibility, R=Repetition, A=Arbitrary Content, ER=Emotional Richness, En=Engagement, H=Hallucinations.}
\label{tab:auto_scores}
\begin{tabular}{lllccccccccc}
\toprule
Run ID & Duration & Run & G & C & Cr & R & A & ER & En & H \\
\midrule
20260420\_054228 & 5 min  & 1 & 5 & 5 & 5 & 5 & 5 & 5 & 5 & 0 \\
20260420\_064300 & 5 min  & 2 & 4 & 5 & 5 & 5 & 5 & 5 & 5 & 1 \\
20260420\_054237 & 10 min & 1 & 5 & 5 & 5 & 5 & 5 & 5 & 5 & 0 \\
20260420\_064335 & 10 min & 2 & 5 & 5 & 5 & 5 & 5 & 5 & 5 & 0 \\
20260420\_054239 & 15 min & 1 & 5 & 5 & 5 & 5 & 5 & 5 & 5 & 0 \\
20260420\_064352 & 15 min & 2 & 3 & 5 & 5 & 5 & 5 & 5 & 5 & 1 \\
\midrule
\textbf{5-min avg}   & & & \textbf{4.50} & \textbf{5.00} & \textbf{5.00} & \textbf{5.00} & \textbf{5.00} & \textbf{5.00} & \textbf{5.00} & \textbf{0.50} \\
\textbf{10-min avg}  & & & \textbf{5.00} & \textbf{5.00} & \textbf{5.00} & \textbf{5.00} & \textbf{5.00} & \textbf{5.00} & \textbf{5.00} & \textbf{0.00} \\
\textbf{15-min avg}  & & & \textbf{4.00} & \textbf{5.00} & \textbf{5.00} & \textbf{5.00} & \textbf{5.00} & \textbf{5.00} & \textbf{5.00} & \textbf{1.00} \\
\textbf{Overall avg} & & & \textbf{4.50} & \textbf{5.00} & \textbf{5.00} & \textbf{5.00} & \textbf{5.00} & \textbf{5.00} & \textbf{5.00} & \textbf{0.33} \\
\bottomrule
\end{tabular}
\end{table}

\noindent\textbf{Key observations.}
\begin{itemize}
  \item \textbf{Grounding is the only variable dimension.} All other six dimensions score a perfect 5.00 across all runs, indicating the pipeline reliably produces coherent, credible, non-repetitive, UI-free, emotionally expressive, and engaging dialogue. Grounding alone varies (range 3--5), reflecting cases where the LLM inferred details not present in the scene graph.
  \item \textbf{Campfire hallucination.} The word \emph{fire} was flagged as a hallucinated entity in two runs (064300 and 064352). This is a systematic artefact of the story prompt framing (``NPC companion at a campfire''), which causes the model to reference fire when no campfire appears in the Minecraft footage. Rephrasing the NPC framing to be environment-neutral would likely eliminate this.
  \item \textbf{Context length and hallucination.} The 10-minute runs achieved perfect Grounding (5.00) with zero hallucinations. The 5-minute and 15-minute runs each produced one hallucination. This may reflect a sweet spot: 5-minute context may be too sparse to anchor the narrative firmly, while 15-minute context may introduce enough entities that the model occasionally infers connections not explicitly stated.
\end{itemize}

\subsection{Human evaluation}

Two raters (Yen and Chris) scored the six Minecraft dialogues independently using the same seven Adams-derived dimensions on a 1--5 scale. Table~\ref{tab:human_scores} reports their ratings alongside the automated judge averages for the same runs.

\begin{table}[h]
\centering
\caption{Human ratings for all six Minecraft runs (scale: 1--5), with automated judge averages for comparison.}
\label{tab:human_scores}
\begin{tabular}{lllccccccc}
\toprule
Run ID & Duration & Rater & G & C & Cr & R & A & ER & En \\
\midrule
20260420\_054228 & 5 min  & Yen   & 4 & 4 & 4 & 5 & 5 & 3 & 4 \\
20260420\_064300 & 5 min  & Chris & 5 & 4 & 3 & 5 & 5 & 3 & 4 \\
20260420\_054237 & 10 min & Yen   & 5 & 4 & 4 & 5 & 5 & 3 & 4 \\
20260420\_064335 & 10 min & Chris & 4 & 4 & 4 & 5 & 5 & 3 & 4 \\
20260420\_054239 & 15 min & Yen   & 5 & 5 & 5 & 5 & 5 & 4 & 4 \\
20260420\_064352 & 15 min & Chris & 5 & 5 & 5 & 5 & 5 & 4 & 4 \\
\midrule
\textbf{Human avg} & & & \textbf{4.67} & \textbf{4.33} & \textbf{4.17} & \textbf{5.00} & \textbf{5.00} & \textbf{3.33} & \textbf{4.00} \\
\textbf{Auto avg}  & & & \textbf{4.50} & \textbf{5.00} & \textbf{5.00} & \textbf{5.00} & \textbf{5.00} & \textbf{5.00} & \textbf{5.00} \\
\bottomrule
\end{tabular}
\end{table}

\noindent\textbf{Key observations.}
\begin{itemize}
  \item \textbf{Strong agreement on Grounding, Repetition, and Arbitrary Content.} Human and automated scores are identical for Repetition and Arbitrary Content (both 5.00) and nearly identical for Grounding (4.67 human vs.\ 4.50 automated), validating these as the most reliable dimensions.
  \item \textbf{Emotional Richness gap.} The largest discrepancy is Emotional Richness: humans averaged 3.33 versus the automated judge's 5.00. The automated judge recognises emotional vocabulary in the text, but cannot assess whether the warmth feels genuine to a human player. This is the primary dimension where the pipeline falls short and where future work should focus.
  \item \textbf{Engagement and Coherence overestimation.} Humans rated Engagement at 4.00 and Coherence at 4.33 versus 5.00 for both from the automated judge, indicating a consistent but moderate overestimation of narrative quality by the automated method.
  \item \textbf{Context length effect.} Both raters gave higher scores to the 15-minute runs (per-rater avg 4.57/5) than to the 5-minute runs (avg 4.14/5), corroborating the finding from the automated scores that richer gameplay context produces more satisfying narratives.
\end{itemize}

\section{Limitations and Future Work}

Frame order is user-supplied (e.g., filename order), not inferred from video timestamps. VLM mistakes propagate into the graph and dialogue; consolidation may over- or under-merge events. The LLM-as-judge uses the same Gemini model family as the generator, contributing to the overestimation of Emotional Richness and Engagement observed in the human study; future work should use an independent judge model or a finer-grained emotional rubric. The campfire framing in the story prompt causes a systematic hallucination of \emph{fire} in environments where no campfire appears; rephrasing the NPC context to be environment-neutral would address this. The human study involved only two raters, each scoring one dialogue per duration group; a larger study with 3--5 raters per sample and inter-rater agreement metrics (e.g., Cohen's $\kappa$) would yield more statistically robust conclusions. The pipeline does not yet integrate live game APIs or native scene graphs from engines; extending to engine telemetry would tighten grounding further.

\bibliographystyle{IEEEbib}
\bibliography{refs}

\end{document}
